% Luc Destombe --- licence CC BY-NC-SA 4.0
% https://creativecommons.org/licenses/by-nc-sa/4.0/deed.fr
% Fichier autonome : compiler avec pdflatex, deux fois (signets, nombre de pages).
\documentclass[11pt,a4paper]{article}
\makeatletter
\newif\ifeleve
\elevefalse
\elevetrue

\newif\iflycee@palatino
\lycee@palatinotrue

\RequirePackage[utf8]{inputenc}
\RequirePackage[T1]{fontenc}
\RequirePackage[french]{babel}
\RequirePackage{etoolbox}

\newcommand{\ShorthandsOff}{\@ifundefined{shorthandoff}{}{\shorthandoff{;:!?}}}
\newcommand{\ShorthandsOn}{\@ifundefined{shorthandon}{}{\shorthandon{;:!?}}}
\ShorthandsOff
\AtBeginDocument{\ShorthandsOn}
\AtBeginEnvironment{tikzpicture}{\ShorthandsOff}

\RequirePackage{geometry}
\RequirePackage{fancyhdr}
\RequirePackage{lastpage}
\RequirePackage{multicol}
\RequirePackage{array}
\RequirePackage{enumitem}
\RequirePackage{xcolor}
\RequirePackage{needspace}

\RequirePackage{amsmath,amssymb}
\RequirePackage{mathpazo}
\DeclareMathSizes{10.95}{10.95}{8}{5.5}
\begingroup\catcode`\_=\active \catcode`\^=\active
\gdef_#1{\sb{\scriptscriptstyle #1}}
\gdef^#1{\sp{\scriptscriptstyle\lycee@moinsepais #1}}
\endgroup
\newcommand\lycee@moins{\mathbin{%
  \setbox\z@\hbox{$\m@th\scriptscriptstyle\lycee@moinsorig$}%
  \dimen@=.55\wd\z@ \advance\dimen@-.5pt
  \hbox to.75\wd\z@{\hss$\m@th\scriptscriptstyle\vcenter{\hbox{%
    \tikz\draw[line width=.5pt,line cap=round](0,0)--(\the\dimen@,0);}}$\hss}}}
\begingroup\lccode`\~=`\- \lowercase{\endgroup\let~}\lycee@moins
\newcommand\lycee@moinsepais{\mathcode`\-="8000 }
\AtBeginDocument{\mathchardef\lycee@moinsorig=\mathcode`\-\relax}
\AtBeginDocument{\catcode`\_=12 \mathcode`\_="8000
  \catcode`\^=12 \mathcode`\^="8000 }
\newcommand\lycee@exposants{%
  \fontdimen13\textfont2=.48\fontdimen6\textfont2
  \fontdimen14\textfont2=.48\fontdimen6\textfont2
  \fontdimen15\textfont2=.48\fontdimen6\textfont2
  \fontdimen13\scriptfont2=.48\fontdimen6\scriptfont2
  \fontdimen14\scriptfont2=.48\fontdimen6\scriptfont2
  \fontdimen15\scriptfont2=.48\fontdimen6\scriptfont2}
\every@math@size\expandafter{\the\every@math@size\lycee@exposants}
\newlength{\retraitcalculs}
\setlength{\retraitcalculs}{2em}

\RequirePackage{tikz}
\usetikzlibrary{arrows.meta,calc,babel,shapes.misc}
\RequirePackage[most]{tcolorbox}
\tcbuselibrary{skins,breakable}

\definecolor{coulDef}{HTML}{1F4E79}
\definecolor{coulProp}{HTML}{7030A0}
\definecolor{coulRem}{HTML}{C55A11}
\definecolor{coulEx}{HTML}{2E7D32}
\definecolor{coulExo}{HTML}{B03A2E}
\definecolor{coulPart}{HTML}{1F4E79}
\definecolor{coulMeth}{HTML}{1B6E4F}
\definecolor{coulCourbe}{HTML}{0B5ED7}
\definecolor{coulCourbeB}{HTML}{C00000}
\definecolor{coulCourbeC}{HTML}{2E7D32}
\definecolor{coulGrille}{HTML}{8C8C8C}
\definecolor{coulRep}{HTML}{1B6E4F}
\definecolor{coulBar}{HTML}{2E7D32}

\newcommand{\R}{\mathbb{R}}
\newcommand{\Rs}{\mathbb{R}^{*}}
\renewcommand{\le}{\leqslant}
\renewcommand{\ge}{\geqslant}
\renewcommand{\approx}{\simeq}
\newcommand{\vect}[1]{\overrightarrow{#1}}
\newcommand{\norme}[1]{\left\lVert #1 \right\rVert}
\newcommand{\intervalle}[2]{\left[#1\,;#2\right]}
\RequirePackage{cancel}
\renewcommand{\CancelColor}{\color{coulExo}}
\newcommand{\fois}[1]{\times{\color{coulExo}#1}}
\newcommand{\barre}[1]{\cancel{#1}}

\tikzset{grille/.style={coulGrille,line width=0.4pt}}
\newcommand{\quadrillage}[4]{\draw[grille,step=1] (#1,#2) grid (#3,#4);}
\tikzset{
  croix/.style={cross out,draw,line width=0.6pt,inner sep=0pt,minimum size=4pt},
  croixdroite/.style={croix,rotate=45,line width=0.8pt},
}
\newcommand{\pt}[3]{\path (#1) node[croix]{} node[#2,font=\small,inner sep=2.5pt]{$#3$};}

\newcommand{\lycee@repere}[4]{%
  \draw[grille,step=1] (#1,#2) grid (#3,#4);
  \draw[-{Stealth[length=2mm]},thick] (#1,0)--({#3+0.4},0) node[below left=-1pt]{$x$};
  \draw[-{Stealth[length=2mm]},thick] (0,#2)--(0,{#4+0.4}) node[below left=-1pt]{$y$};
  \node[below left,font=\scriptsize,inner sep=1.5pt] at (0,0) {$0$};}
\newcommand{\lycee@gradx}[1]{\foreach \k/\l in {#1}{%
  \draw (\k,0.07)--(\k,-0.07) node[below,font=\scriptsize,inner sep=1.5pt]{$\l$};}}
\newcommand{\lycee@grady}[1]{\foreach \k/\l in {#1}{%
  \draw (0.07,\k)--(-0.07,\k) node[left,font=\scriptsize,inner sep=1.5pt]{$\l$};}}
\newcommand{\lycee@intff}[2]{\left[#1\,;#2\right]}
\newcommand{\lycee@intfo}[2]{\left[#1\,;#2\right[}
\newcommand{\lycee@intof}[2]{\left]#1\,;#2\right]}
\newcommand{\lycee@intoo}[2]{\left]#1\,;#2\right[}
\newcommand{\lycee@ens}[1]{\left\{#1\right\}}
\AtBeginDocument{%
  \providecommand{\repere}{\lycee@repere}%
  \providecommand{\gradx}{\lycee@gradx}%
  \providecommand{\grady}{\lycee@grady}%
  \providecommand{\Cf}{\mathcal{C}\sb{\scriptscriptstyle f}}%
  \providecommand{\Cg}{\mathcal{C}\sb{\scriptscriptstyle g}}%
  \providecommand{\intff}{\lycee@intff}%
  \providecommand{\intfo}{\lycee@intfo}%
  \providecommand{\intof}{\lycee@intof}%
  \providecommand{\intoo}{\lycee@intoo}%
  \providecommand{\ens}{\lycee@ens}}

\newlist{questions}{enumerate}{3}
\setlist[questions,1]{label=\arabic*\textdegree),leftmargin=*,itemsep=2pt,topsep=3pt}
\setlist[questions,2]{label=\alph*),leftmargin=*,itemsep=1pt,topsep=2pt}
\setlist[questions,3]{label=\roman*),leftmargin=*,itemsep=1pt,topsep=1pt}
\setlist[itemize]{label=\textbullet,leftmargin=*,itemsep=2pt,topsep=3pt}

\newcommand{\besoinplace}[1]{\par\penalty\@M\begingroup
  \dimen@=#1\relax\dimen@ii=\pagegoal\advance\dimen@ii-\pagetotal
  \ifdim\dimen@>\dimen@ii\ifdim\dimen@ii>\z@\vfil\fi\break\fi\endgroup}

\newcommand{\lycee@pied}{\thepage/\pageref{LastPage}}
\newcommand{\entete}[2]{%
  \pagestyle{fancy}\fancyhf{}%
  \fancyhead[L]{\footnotesize\color{coulPart}\textbf{#1}}%
  \fancyhead[R]{\footnotesize\color{coulPart}\textbf{#2}}%
  \fancyfoot[C]{\footnotesize\lycee@pied}%
  \renewcommand{\headrulewidth}{0.4pt}%
  \renewcommand{\headrule}{\hbox to\headwidth{%
    \color{coulPart!40}\leaders\hrule height \headrulewidth\hfill}}}

\RequirePackage{ccicons}
\newcommand{\lycee@licence}{{\scriptsize\color{gray}%
  \copyright\ Luc Destombe --- \ccbyncsa\ CC BY-NC-SA 4.0}}
\newif\iflycee@licence \lycee@licencetrue
\newcommand{\sanslicence}{\lycee@licencefalse}
\AtBeginDocument{\iflycee@licence\AddToHookNext{shipout/foreground}{%
  \put(0,0){\rlap{\kern\dimexpr1in+\hoffset+\oddsidemargin+\textwidth\relax
    \llap{\raisebox{-\dimexpr1in+\voffset+\topmargin+\headheight+\headsep
      +\textheight+\footskip\relax}{\lycee@licence}}}}}\fi}

\newcommand{\formulesserrees}{%
  \setlength{\abovedisplayskip}{3pt}\setlength{\belowdisplayskip}{3pt}%
  \setlength{\abovedisplayshortskip}{2pt}\setlength{\belowdisplayshortskip}{3pt}}

\setlength{\parindent}{0pt}

\RequirePackage[implicit=false,hidelinks,pdfencoding=auto,psdextra,bookmarksopen,
  bookmarksopenlevel=1,pdfcreator={},pdfproducer={}]{hyperref}
\RequirePackage{bookmark}
\hypersetup{pdfauthor={Luc Destombe},pdfkeywords={CC BY-NC-SA 4.0}}
\newcounter{lycee@signet}
\newcommand{\signet}[3][]{\stepcounter{lycee@signet}%
  \edef\lycee@dest{\if\relax\detokenize{#1}\relax
    signet.\arabic{lycee@signet}\else#1\fi}%
  \hypertarget{\lycee@dest}{}\bookmark[level=#2,dest=\lycee@dest]{#3}}
\pdfstringdefDisableCommands{%
  \def\R{R}\def\vect#1{#1}\def\Cf{Cf}\def\Cg{Cg}\def\fois#1{×#1}%
  \def\barre#1{#1}\def\intervalle#1#2{[#1;#2]}\def\intff#1#2{[#1;#2]}%
  \def\textsuperscript#1{#1}\def\dfrac#1#2{#1/#2}\def\frac#1#2{#1/#2}%
  \def\vec#1{#1⃗}}%
\geometry{top=2cm,bottom=1cm,left=1cm,right=1cm,headheight=14pt,footskip=0.6cm}

\RequirePackage{pgfplots}
\pgfplotsset{compat=1.18}
\RequirePackage{tkz-tab}

\newcounter{partiecours}
\renewcommand{\thepartiecours}{\Roman{partiecours}}
\newcommand{\partiecours}[1]{%
  \refstepcounter{partiecours}%
  \Needspace*{8\baselineskip}%
  \bigskip
  \noindent\begin{tcolorbox}[enhanced,colback=coulPart,colframe=coulPart,
      boxrule=0pt,arc=2pt,left=8pt,right=8pt,top=4pt,bottom=4pt,
      before skip=6pt,after skip=10pt]
    \color{white}\bfseries\large \thepartiecours{} --- #1%
    \signet[partie-\arabic{partiecours}]{1}{\thepartiecours{} --- #1}
  \end{tcolorbox}%
  \addcontentsline{toc}{section}{\thepartiecours{} --- #1}%
}

\newtcolorbox{boitecours}[3][]{%
  enhanced,breakable,
  colback=white,colframe=#2,coltitle=white,
  fonttitle=\bfseries,
  attach boxed title to top left={xshift=8pt,yshift=-\tcboxedtitleheight/2},
  boxed title style={colback=#2,arc=2pt,boxrule=0pt},
  boxrule=0.9pt,arc=2pt,
  left=8pt,right=8pt,top=8pt,bottom=6pt,
  title={#3},#1}

\newcounter{methode}
\newenvironment{definitionbox}[1][Définition]%
  {\begin{boitecours}[phantom={\signet{2}{#1}}]{coulDef}{#1}}{\end{boitecours}}
\newenvironment{proprietebox}[1][Propriété]%
  {\begin{boitecours}[phantom={\signet{2}{#1}}]{coulProp}{#1}}{\end{boitecours}}
\newenvironment{methodebox}[1][Méthode]{\stepcounter{methode}%
  \begin{boitecours}[phantom={\signet[methode-\arabic{methode}]{2}{#1}}]{coulMeth}{#1}}%
  {\end{boitecours}}

\newcommand{\etape}[1]{\par\smallskip\textbf{\color{coulMeth}#1}\par\nobreak\smallskip}
\newcommand{\enonce}[1]{\emph{#1}\par\smallskip}
\newcommand{\reponse}[1]{\par\smallskip
  {\footnotesize\itshape\color{black!55}Réponses : #1}}
\newcommand{\demo}[1]{\textbf{\color{coulMeth}Démonstration.} #1}
\newcommand{\afaire}[1]{%
  \par\smallskip
  \noindent{\small\color{coulExo}$\blacktriangleright$\ \textbf{À faire :}
    fiche d'exercices du chapitre, #1.}\par\smallskip}

\newenvironment{calculs}{\setlength{\jot}{5pt}\@fleqntrue\@mathmargin\retraitcalculs
  \setlength{\abovedisplayskip}{4pt}\setlength{\belowdisplayskip}{4pt}%
  \setlength{\abovedisplayshortskip}{4pt}\setlength{\belowdisplayshortskip}{4pt}%
  \csname alignat*\endcsname{2}}%
  {\csname endalignat*\endcsname}
\newcommand{\rem}[1]{\text{\small\itshape\color{black!60}#1}}
\newcommand{\explique}[2]{\par\smallskip\noindent
  \begin{minipage}[t]{0.57\linewidth}\raggedright #1\end{minipage}\hfill
  \begin{minipage}[t]{0.40\linewidth}\raggedright\leavevmode
    {\small\itshape\color{black!60}#2}\end{minipage}\par}
\newcommand{\cas}[1]{\par\smallskip\noindent\textbf{\color{coulExo}#1}\nobreak}
\newcommand{\calc}[1]{\par\smallskip\textbf{\color{coulExo}#1}\par\nobreak}

\newtcolorbox{remarquebox}[1][Remarque]{%
  enhanced,breakable,colback=white,colframe=coulRem,
  boxrule=0pt,leftrule=2.5pt,arc=0pt,outer arc=0pt,
  left=8pt,right=6pt,top=5pt,bottom=5pt,
  before upper={\textbf{\color{coulRem}#1}\par\smallskip}}

\newtcolorbox{exemplebox}[1][Exemple]{%
  enhanced,breakable,colback=white,colframe=coulEx,
  boxrule=0pt,leftrule=2.5pt,arc=0pt,outer arc=0pt,
  left=8pt,right=6pt,top=5pt,bottom=5pt,
  before upper={\textbf{\color{coulEx}#1}\par\smallskip}}

\pgfplotsset{
  repere/.style={
    axis lines=middle,
    axis line style={-{Stealth[length=2.4mm]},thick},
    every axis x label/.style={at={(ticklabel* cs:1.0)},anchor=west},
    every axis y label/.style={at={(ticklabel* cs:1.0)},anchor=south},
    label style={font=\footnotesize},
    tick label style={font=\scriptsize},
    grid=both,
    major grid style={grille},
    minor tick num=0,
    tick align=inside,
    clip mode=individual,
  },
}
\tikzset{
  courbe/.style={coulCourbe,line width=1.1pt,smooth},
  courbeB/.style={coulCourbeB,line width=1.1pt,smooth},
  courbeC/.style={coulCourbeC,line width=1.1pt,smooth},
}

\newlist{etapes}{enumerate}{1}
\setlist[etapes]{label=\textbf{\arabic*.},leftmargin=*,itemsep=1pt,topsep=2pt}

\newlist{expressions}{enumerate}{1}
\setlist[expressions]{label={\Alph*\ $=$},leftmargin=*,itemsep=3pt,topsep=3pt}

\newcounter{exo}
\renewcommand{\theexo}{\ifnum\value{exo}<10 0\fi\arabic{exo}}
\newtcolorbox{exercice}[1][]{%
  enhanced,breakable,
  colback=white,colframe=coulExo,
  boxrule=0.9pt,arc=2pt,
  left=8pt,right=8pt,top=8pt,bottom=6pt,
  coltitle=white,fonttitle=\bfseries,
  attach boxed title to top left={xshift=8pt,yshift=-\tcboxedtitleheight/2},
  boxed title style={colback=coulExo,arc=2pt,boxrule=0pt},
  phantom={\signet{2}{Exercice \arabic{exo}}},
  title={Exercice \theexo\ifx\relax#1\relax\else{} --- #1\fi}}
\newenvironment{exo}[1][\relax]{\refstepcounter{exo}\begin{exercice}[#1]}{\end{exercice}}

  \RequirePackage{comment}
  \excludecomment{correction}
  \renewcommand{\reponse}[1]{}
  \newcommand{\autableau}{\hfill{\footnotesize\itshape\color{black!45}(au tableau)}}
  \newcommand{\etapeexemples}{%
    \par\smallskip\textbf{\color{coulMeth}Exemples}\autableau\par\nobreak\smallskip}
  \newcommand{\etapetableau}[1]{%
    \par\smallskip\textbf{\color{coulMeth}#1}\autableau\par\nobreak\smallskip}
  \newtcolorbox{exempletableau}[1][Exemple]{%
    enhanced,breakable,colback=white,colframe=coulEx,
    boxrule=0pt,leftrule=2.5pt,arc=0pt,outer arc=0pt,
    left=8pt,right=6pt,top=4pt,bottom=4pt,
    before skip=5pt,after skip=5pt,
    before upper={\textbf{\color{coulEx}#1}\autableau\par\smallskip}}

\RequirePackage{comment}
\excludecomment{correctionavj}

\newcommand{\coursEntete}{}
\newcommand{\coursNiveau}{}
\newcommand{\coursTitre}{}
\newcommand{\coursSurTitre}{}
\newcommand{\coursSousTitre}{}

\newcommand{\enteteCours}[1]{\renewcommand{\coursEntete}{#1}}
\newcommand{\chapitrecours}[2]{\enteteCours{Chapitre #1 --- #2}}
\newcommand{\niveaucours}[1]{\renewcommand{\coursNiveau}{#1}}
\newcommand{\titrecours}[3][]{%
  \renewcommand{\coursTitre}{#2}%
  \renewcommand{\coursSurTitre}{#1}%
  \renewcommand{\coursSousTitre}{#3}}

\pagestyle{fancy}
\fancyhf{}
\fancyhead[L]{\footnotesize\color{coulPart}\textbf{\coursEntete}}
\fancyhead[R]{\footnotesize\color{coulPart}\textbf{\coursNiveau}}
\fancyfoot[C]{\footnotesize\thepage}
\renewcommand{\headrulewidth}{0.4pt}
\renewcommand{\headrule}{\hbox to\headwidth{\color{coulPart!40}\leaders\hrule height \headrulewidth\hfill}}

\setlength{\parskip}{2pt}

\AtBeginDocument{%
  \begin{center}
    \begin{tcolorbox}[enhanced,width=0.72\linewidth,colback=white,colframe=coulPart,
        boxrule=1.6pt,arc=3pt,halign=center,top=10pt,bottom=10pt]
      \ifx\coursSurTitre\empty
        {\Huge\bfseries\color{coulPart} \coursTitre}\\[4pt]
      \else
        {\Huge\bfseries\color{coulPart} \coursTitre}\\[3pt]
        {\Large\bfseries\color{coulPart} \coursSurTitre}\\[5pt]
      \fi
      {\normalsize \coursSousTitre}%
      \\[2pt]{\small\itshape Version élève}
    \end{tcolorbox}
  \end{center}%
}
\makeatother

\chapitrecours{6}{Probabilités conditionnelles}
\niveaucours{1\textsuperscript{re} spécialité mathématiques}
\titrecours{PROBABILITÉS CONDITIONNELLES}{Cours --- Première spécialité}

\begin{document}

\tikzset{evt/.style={inner sep=2pt},
  proba/.style={font=\footnotesize,fill=white,inner sep=1pt},
  nom/.style={font=\footnotesize,inner sep=1.5pt,sloped},
  arbA/.style={},arbnA/.style={},arbAB/.style={},arbAnB/.style={},arbnAB/.style={},
  arbnAnB/.style={},chemin/.style={coulExo,line width=1.4pt}}
\newcommand{\arbcontraires}[2]{\def\arbnA{#1}\def\arbnB{#2}}
\newcommand{\arbfeuilles}[4]{\def\arbfAB{#1}\def\arbfAnB{#2}\def\arbfnAB{#3}\def\arbfnAnB{#4}}
\newcommand{\arbsansfeuilles}{\def\arbpasfeuilles{}}
\newcommand{\arbvaleurs}{\def\arbseules{}\arbsansfeuilles}
\newcommand{\arblab}[2]{\ifx\arbseules\undefined
  $#1\if\relax\detokenize{#2}\relax\else=#2\fi$\else$#2$\fi}
\newcommand{\arbfeuille}[2]{%
  \node[right,font=\footnotesize] at (#1.east) {$#2\ifx\arbtmp\empty\else=\arbtmp\fi$};}
\newcommand{\arbre}[8]{%
  \ifx\arbnA\undefined\def\arbNA{\overline{#1}}\else\let\arbNA\arbnA\fi
  \ifx\arbnB\undefined\def\arbNB{\overline{#2}}\else\let\arbNB\arbnB\fi
  \node (O) at (0,0) {};
  \node[evt] (A) at (3.4,1.2) {$#1$};
  \node[evt] (nA) at (3.4,-1.2) {$\arbNA$};
  \node[evt] (AB) at (7.2,1.9) {$#2$};
  \node[evt] (AnB) at (7.2,0.5) {$\arbNB$};
  \node[evt] (nAB) at (7.2,-0.5) {$#2$};
  \node[evt] (nAnB) at (7.2,-1.9) {$\arbNB$};
  \draw[arbA] (O.center) -- node[nom,above]{\arblab{P(#1)}{#3}} (A);
  \draw[arbnA] (O.center) -- node[nom,below]{\arblab{P(\arbNA)}{#4}} (nA);
  \draw[arbAB] (A) -- node[nom,above]{\arblab{P_{#1}(#2)}{#5}} (AB);
  \draw[arbAnB] (A) -- node[nom,below]{\arblab{P_{#1}(\arbNB)}{#6}} (AnB);
  \draw[arbnAB] (nA) -- node[nom,above]{\arblab{P_{\arbNA}(#2)}{#7}} (nAB);
  \draw[arbnAnB] (nA) -- node[nom,below]{\arblab{P_{\arbNA}(\arbNB)}{#8}} (nAnB);
  \ifx\arbpasfeuilles\undefined
    \ifx\arbfAB\undefined\arbfeuilles{}{}{}{}\fi
    \let\arbtmp\arbfAB\arbfeuille{AB}{P(#1\cap #2)}
    \let\arbtmp\arbfAnB\arbfeuille{AnB}{P(#1\cap\arbNB)}
    \let\arbtmp\arbfnAB\arbfeuille{nAB}{P(\arbNA\cap #2)}
    \let\arbtmp\arbfnAnB\arbfeuille{nAnB}{P(\arbNA\cap\arbNB)}
  \fi}
\newcommand{\feuille}[2]{\node[right,font=\footnotesize] at (#1.east) {#2};}
\newcommand{\arbremachines}[3]{%
  \node (O) at (0,0) {};
  \foreach \p/\q/\r [count=\i] in {#3}{%
    \pgfmathsetmacro{\y}{4.6-2.3*\i}
    \node[evt] (M\i) at (3.4,\y) {$#1_{\i}$};
    \ifnum\i=3 \tikzset{arbpos/.style={below}}\else\tikzset{arbpos/.style={above}}\fi
    \draw (O.center) -- node[nom,arbpos]{$P(#1_{\i})=\p$} (M\i);
    \node[evt] (D\i) at (7.2,\y+0.6) {$#2$};
    \node[evt] (nD\i) at (7.2,\y-0.6) {$\overline{#2}$};
    \draw (M\i) -- node[nom,above]{$P_{#1_{\i}}(#2)=\q$} (D\i);
    \draw (M\i) -- node[nom,below]{$P_{#1_{\i}}(\overline{#2})={\color{coulExo}\r}$} (nD\i);
    \node[right,font=\footnotesize] at (D\i.east) {$P(#1_{\i}\cap #2)$};
    \node[right,font=\footnotesize] at (nD\i.east) {$P(#1_{\i}\cap\overline{#2})$};}}
\newcommand{\arbretrois}[2]{%
  \node (O) at (0,0) {};
  \foreach \a/\ya/\na/\pa in {1/1.8/S/#1,0/-1.8/\overline{S}/#2}{%
    \node[evt] (N\a) at (2.2,\ya) {$\na$};
    \draw (O.center) -- node[proba]{$\pa$} (N\a);
    \foreach \b/\db/\nb/\pb in {1/0.9/S/#1,0/-0.9/\overline{S}/#2}{%
      \pgfmathsetmacro{\yb}{\ya+\db}
      \node[evt] (N\a\b) at (4.4,\yb) {$\nb$};
      \draw (N\a) -- node[proba]{$\pb$} (N\a\b);
      \foreach \c/\dc/\nc/\pc in {1/0.45/S/#1,0/-0.45/\overline{S}/#2}{%
        \pgfmathsetmacro{\yc}{\yb+\dc}
        \node[evt] (N\a\b\c) at (7,\yc) {$\nc$};
        \draw (N\a\b) -- node[proba]{$\pc$} (N\a\b\c);}}}}
\newcommand{\vennfond}{%
  \draw (-0.85,-0.75) rectangle (1.65,0.75);
  \node[font=\footnotesize] at (-0.65,0.53) {$\Omega$};}

\begin{remarquebox}[Comment utiliser ce cours]
Chaque partie commence par ce qu'il faut \textbf{savoir} (définitions, théorèmes), puis
vient ce qu'il faut \textbf{savoir faire} : une méthode, avec des
\textbf{exemples} faits en classe, puis des exercices
\og\textbf{À vous de jouer}\fg{} à chercher seul.

\smallskip
Sur cette feuille, seuls les \textbf{énoncés} des exemples figurent : la résolution,
signalée par la mention \emph{(au tableau)}, est à rédiger dans le cahier.

\end{remarquebox}

\partiecours{Vocabulaire des probabilités (rappels de seconde)}

\begin{definitionbox}[Définition 1 --- Expérience aléatoire, issues, événements]
\begin{itemize}
  \item Une \textbf{expérience aléatoire} a plusieurs résultats possibles, appelés
    \textbf{issues}, et on ne peut pas prévoir lequel sortira.
  \item L'ensemble de toutes les issues est l'\textbf{univers}, noté $\Omega$.
    Pour un dé à six faces : $\Omega=\ens{1\,;2\,;3\,;4\,;5\,;6}$.
  \item Un \textbf{événement} est un ensemble d'issues. Pour le dé, \og obtenir un
    nombre pair\fg{} est l'événement $A=\ens{2\,;4\,;6}$.
  \item $\Omega$ est l'événement \textbf{certain} ; l'ensemble vide $\varnothing$ est
    l'événement \textbf{impossible}.
\end{itemize}
\end{definitionbox}

\begin{definitionbox}[Définition 2 --- Contraire, intersection, réunion]
Soit $A$ et $B$ deux événements.
\begin{itemize}
  \item Le \textbf{contraire} de $A$, noté $\overline{A}$ (\og non $A$\fg), contient
    les issues qui ne sont pas dans $A$.
  \item L'\textbf{intersection} $A\cap B$ (\og $A$ \textbf{et} $B$\fg) contient les
    issues qui sont à la fois dans $A$ et dans $B$.
  \item La \textbf{réunion} $A\cup B$ (\og $A$ \textbf{ou} $B$\fg) contient les issues
    qui sont dans $A$, dans $B$ ou dans les deux : au moins l'un des deux est réalisé.
  \item $A$ et $B$ sont \textbf{incompatibles} quand $A\cap B=\varnothing$ : ils ne
    peuvent pas se produire ensemble.
\end{itemize}
\begin{center}
\begin{tikzpicture}
  \vennfond
  \fill[coulMeth!30] (-0.85,-0.75) rectangle (1.65,0.75);
  \fill[white] (0.4,0) circle (0.55);
  \draw (-0.85,-0.75) rectangle (1.65,0.75) (0.4,0) circle (0.55);
  \node[font=\footnotesize] at (0.4,0) {$A$};
  \node[font=\footnotesize] at (1.3,-0.4) {$\overline{A}$};
  \node[below,font=\footnotesize] at (0.4,-0.75) {contraire};
\end{tikzpicture}\hspace{0.8cm}
\begin{tikzpicture}
  \vennfond
  \begin{scope}\clip (0,0) circle (0.55); \fill[coulMeth!30] (0.8,0) circle (0.55);\end{scope}
  \draw (0,0) circle (0.55) (0.8,0) circle (0.55);
  \node[font=\footnotesize] at (-0.25,0) {$A$};
  \node[font=\footnotesize] at (1.05,0) {$B$};
  \node[font=\footnotesize] (I) at (0.4,-1.05) {$A\cap B$};
  \draw[->,>=stealth] (I.north) -- (0.4,0);
\end{tikzpicture}\hspace{0.8cm}
\begin{tikzpicture}
  \vennfond
  \fill[coulMeth!30] (0,0) circle (0.55) (0.8,0) circle (0.55);
  \draw (0,0) circle (0.55) (0.8,0) circle (0.55);
  \node[font=\footnotesize] at (-0.25,0.15) {$A$};
  \node[font=\footnotesize] at (1.05,0.15) {$B$};
  \node[font=\footnotesize] (U) at (0.4,-1.05) {$A\cup B$};
  \draw[->,>=stealth] (U.north) -- (-0.15,-0.2);
  \draw[->,>=stealth] (U.north) -- (0.4,-0.1);
  \draw[->,>=stealth] (U.north) -- (0.95,-0.2);
\end{tikzpicture}\hspace{0.8cm}
\begin{tikzpicture}
  \vennfond
  \draw (-0.15,0) circle (0.45) (0.95,0) circle (0.45);
  \node[font=\footnotesize] at (-0.15,0) {$A$};
  \node[font=\footnotesize] at (0.95,0) {$B$};
  \node[below,font=\footnotesize] at (0.4,-0.75) {incompatibles : $A\cap B=\varnothing$};
\end{tikzpicture}
\end{center}
\end{definitionbox}

\begin{proprietebox}[Propriété 1 --- Calculer une probabilité]
\begin{itemize}
  \item Pour tout événement $A$ : $0\le P(A)\le 1$, avec $P(\Omega)=1$ et
    $P(\varnothing)=0$. La probabilité de $A$ est la somme des probabilités de ses issues.
  \item \textbf{Équiprobabilité} : quand toutes les issues ont la même probabilité
    (tirage \og au hasard\fg, dé équilibré),
    \[
      P(A)=\frac{\text{nombre d'issues de }A}{\text{nombre d'issues de }\Omega}.
    \]
  \item $P(\overline{A})=1-P(A)$.
  \item $P(A\cup B)=P(A)+P(B)-P(A\cap B)$ : on retire $P(A\cap B)$, compté deux fois.
\end{itemize}
\end{proprietebox}

\begin{methodebox}[Méthode 1 --- Décrire des événements et calculer leurs probabilités]
\etapeexemples
Une application de quiz choisit au hasard une question parmi $20$, numérotées de $1$
à $20$. On note $A$ : \og le numéro est pair\fg{} et $B$ : \og le numéro est un multiple
de $5$\fg.
\begin{questions}[label=\alph*)]
  \item Écrire $A$, $B$, $A\cap B$ et $A\cup B$ sous forme d'ensembles.
  \item Calculer $P(A)$, $P(B)$ et $P(A\cap B)$.
  \item Calculer $P(A\cup B)$ avec la formule, puis vérifier en comptant.
  \item Décrire $\overline{A}$ par une phrase et calculer $P(\overline{A})$.
\end{questions}

\etape{À vous de jouer}
Au départ d'une course, on tire au hasard un dossard parmi $30$, numérotés de $1$ à
$30$. On note $A$ : \og le numéro est un multiple de $3$\fg{} et $B$ : \og le numéro
est un multiple de $5$\fg.
\begin{questions}[label=\alph*)]
  \item Écrire $A$, $B$, $A\cap B$ et $A\cup B$ sous forme d'ensembles.
  \item Calculer $P(A)$, $P(B)$ et $P(A\cap B)$.
  \item Calculer $P(A\cup B)$ avec la formule, puis vérifier en comptant.
  \item Décrire $\overline{A}$ par une phrase et calculer $P(\overline{A})$.
\end{questions}

\end{methodebox}

\begin{methodebox}[Méthode 2 --- Calculer avec les formules, sans compter les issues]
\etapeexemples
On choisit au hasard un client d'un site de vente en ligne. On note $A$ : \og le
client a acheté un livre\fg{} et $B$ : \og le client a acheté un jeu\fg. On sait que
$P(A)=0{,}35$, $P(B)=0{,}25$ et $P(A\cap B)=0{,}1$.
\begin{questions}[label=\alph*)]
  \item Décrire par une phrase les événements $A\cap B$ et $A\cup B$.
  \item Calculer $P(A\cup B)$.
  \item Calculer la probabilité que le client n'ait acheté ni livre ni jeu.
\end{questions}

\etape{À vous de jouer}
On choisit au hasard un utilisateur d'un téléphone. On note $A$ : \og il est abonné à
une plateforme de musique\fg{} et $B$ : \og il est abonné à une plateforme de vidéo\fg.
On sait que $P(A)=0{,}6$, $P(B)=0{,}3$ et $P(A\cap B)=0{,}2$.
\begin{questions}[label=\alph*)]
  \item Décrire par une phrase les événements $A\cap B$ et $A\cup B$.
  \item Calculer $P(A\cup B)$.
  \item Calculer la probabilité que l'utilisateur ne soit abonné à aucune des deux.
\end{questions}

\end{methodebox}

\afaire{exercices 1 à 4}

\partiecours{Probabilités conditionnelles et tableaux}

\begin{definitionbox}[Définition 3 --- Probabilité de $B$ sachant $A$]
\begin{minipage}[c]{0.62\linewidth}\raggedright
Soit $A$ un événement de probabilité non nulle. La \textbf{probabilité de $B$ sachant
$A$}, notée $P_{A}(B)$, est la probabilité que $B$ se réalise quand on sait que $A$
est réalisé : on ne regarde plus que les issues de $A$.

\smallskip
En situation d'équiprobabilité (tableau d'effectifs) :
\[
  P_{A}(B)=\frac{\text{effectif de }A\cap B}{\text{effectif de }A}.
\]
\end{minipage}\hfill
\begin{minipage}[c]{0.34\linewidth}\centering
\begin{tikzpicture}
  \vennfond
  \begin{scope}\clip (0,0) circle (0.55); \fill[coulMeth!30] (0.8,0) circle (0.55);\end{scope}
  \draw[coulMeth,line width=1.4pt] (0,0) circle (0.55);
  \draw (0.8,0) circle (0.55);
  \node[font=\footnotesize] at (-0.25,0) {$A$};
  \node[font=\footnotesize] at (1.05,0) {$B$};
  \node[below,font=\footnotesize,align=center] at (0.4,-0.75) {$A$ devient\\ l'univers};
\end{tikzpicture}
\end{minipage}
\end{definitionbox}

\begin{remarquebox}[Lire l'énoncé : « et » ou « parmi » ?]
\begin{itemize}
  \item \og $18\,\%$ des élèves sont en seconde \textbf{et} viennent en bus\fg{} : on
    parle de tous les élèves, c'est $P(S\cap T)=0{,}18$.
  \item \og \textbf{Parmi} les élèves de seconde, $60\,\%$ viennent en bus\fg{} (ou
    \og \textbf{sachant} qu'il est en seconde\fg, \og \textbf{si} l'élève est en
    seconde\fg) : on ne parle que des élèves de seconde, c'est $P_{S}(T)=0{,}6$.
  \item En général, $P_{A}(B)$, $P_{B}(A)$ et $P(A\cap B)$ sont trois nombres différents.
\end{itemize}
\end{remarquebox}

\begin{methodebox}[Méthode 3 --- Calculer une probabilité conditionnelle dans un tableau]
\etapeexemples
On a demandé aux $400$ élèves d'un lycée s'ils viennent en bus.
\begin{center}
\renewcommand{\arraystretch}{1.3}
\begin{tabular}{|l|c|c|c|}
\hline
 & Bus & Autre moyen & Total\\\hline
Seconde & $120$ & $80$ & $200$\\\hline
Première & $90$ & $110$ & $200$\\\hline
Total & $210$ & $190$ & $400$\\\hline
\end{tabular}
\end{center}
On choisit un élève au hasard. On note $S$ : \og l'élève est en seconde\fg{} et
$T$ : \og l'élève vient en bus\fg.
\begin{questions}[label=\alph*)]
  \item Calculer $P(S)$, $P(T)$ et $P(S\cap T)$.
  \item L'élève choisi est en seconde. Calculer la probabilité qu'il vienne en bus.
  \item L'élève choisi vient en bus. Calculer la probabilité qu'il soit en seconde.
\end{questions}

\etape{À vous de jouer}
Un cinéma a compté ses $250$ spectateurs d'une séance.
\begin{center}
\renewcommand{\arraystretch}{1.3}
\begin{tabular}{|l|c|c|c|}
\hline
 & Pop-corn & Sans pop-corn & Total\\\hline
Enfant & $60$ & $40$ & $100$\\\hline
Adulte & $45$ & $105$ & $150$\\\hline
Total & $105$ & $145$ & $250$\\\hline
\end{tabular}
\end{center}
On choisit un spectateur au hasard. On note $E$ : \og c'est un enfant\fg{} et
$C$ : \og il a acheté du pop-corn\fg.
\begin{questions}[label=\alph*)]
  \item Calculer $P(E)$, $P(C)$ et $P(E\cap C)$.
  \item Le spectateur choisi est un enfant. Calculer la probabilité qu'il ait acheté du
    pop-corn.
  \item Le spectateur choisi a acheté du pop-corn. Calculer la probabilité que ce soit
    un enfant.
\end{questions}

\end{methodebox}

\afaire{exercices 5 à 8}

\partiecours{Formule des probabilités conditionnelles}

\begin{proprietebox}[Propriété 2 --- Formule des probabilités conditionnelles]
Soit $A$ un événement tel que $P(A)\neq 0$ et $B$ un événement. Alors
\[
  P_{A}(B)=\frac{P(A\cap B)}{P(A)}
  \qquad\text{et donc}\qquad
  P(A\cap B)=P(A)\times P_{A}(B).
\]
De plus, $P_{A}(\overline{B})=1-P_{A}(B)$.
\end{proprietebox}

\demo{dans un tableau d'effectifs, on note $n(A)$ le nombre d'issues de $A$ et $N$ le
nombre total d'issues.}
\begin{calculs}
P_{A}(B) &= \dfrac{n(A\cap B)}{n(A)} &\qquad& \rem{on compte dans $A$ seulement}\\
P_{A}(B) &= \dfrac{n(A\cap B)\div N}{n(A)\div N} && \rem{on divise en haut et en bas
par $N$}\\
P_{A}(B) &= \dfrac{P(A\cap B)}{P(A)} && \rem{chaque quotient est une probabilité}
\end{calculs}

\begin{methodebox}[Méthode 4 --- Traduire un énoncé et utiliser la formule]
\etapeexemples
On choisit au hasard un abonné d'une plateforme de vidéos. On note $F$ : \og il a
l'offre famille\fg{} et $T$ : \og il regarde sur une télévision\fg.
\begin{itemize}
  \item $40\,\%$ des abonnés ont l'offre famille ;
  \item $12\,\%$ des abonnés ont l'offre famille et regardent sur une télévision ;
  \item parmi les abonnés qui n'ont pas l'offre famille, $25\,\%$ regardent sur une
    télévision.
\end{itemize}
\begin{questions}[label=\alph*)]
  \item Traduire chaque donnée par une probabilité.
  \item Calculer $P_{F}(T)$ et interpréter le résultat.
  \item Calculer la probabilité que l'abonné n'ait pas l'offre famille et regarde sur
    une télévision.
\end{questions}

\etape{À vous de jouer}
On choisit au hasard un client d'une salle de sport. On note $A$ : \og il a un
abonnement annuel\fg{} et $C$ : \og il prend des cours collectifs\fg.
\begin{itemize}
  \item $60\,\%$ des clients ont un abonnement annuel ;
  \item $15\,\%$ des clients ont un abonnement annuel et prennent des cours collectifs ;
  \item parmi les clients sans abonnement annuel, $30\,\%$ prennent des cours collectifs.
\end{itemize}
\begin{questions}[label=\alph*)]
  \item Traduire chaque donnée par une probabilité.
  \item Calculer $P_{A}(C)$ et interpréter le résultat.
  \item Calculer la probabilité que le client n'ait pas d'abonnement annuel et prenne
    des cours collectifs.
\end{questions}

\end{methodebox}

\afaire{exercices 9 à 12}

\partiecours{Arbres pondérés}

\begin{proprietebox}[Propriété 3 --- Règles de l'arbre pondéré]
\begin{minipage}[c]{0.44\linewidth}\raggedright
\begin{itemize}
  \item Au premier niveau : $P(A)$ et $P(\overline{A})$.
  \item Au second niveau : des probabilités \textbf{conditionnelles} (sachant
    l'événement d'où part la branche).
  \item \textbf{Règle du nœud} : les branches issues d'un même nœud ont une somme
    égale à $1$.
  \item \textbf{Règle du chemin} : la probabilité d'un chemin est le
    \textbf{produit} des probabilités de ses branches ; c'est la probabilité de
    l'intersection, par exemple $P(A\cap B)=P(A)\times P_{A}(B)$.
\end{itemize}
\end{minipage}\hfill
\begin{minipage}[c]{0.54\linewidth}\centering
\begin{tikzpicture}
  \arbre{A}{B}{}{}{}{}{}{}
\end{tikzpicture}
\end{minipage}
\end{proprietebox}

\begin{methodebox}[Méthode 5 --- Compléter un arbre et lire une intersection]
\etapeexemples
\begin{minipage}[c]{0.44\linewidth}\raggedright
On donne l'arbre pondéré ci-contre.
\begin{questions}[label=\alph*)]
  \item Écrire chaque nombre donné sous la forme d'une probabilité.
  \item Compléter l'arbre.
  \item Calculer $P(A\cap\overline{B})$ et $P(\overline{A}\cap B)$.
\end{questions}
\end{minipage}\hfill
\begin{minipage}[c]{0.54\linewidth}\centering
\begin{tikzpicture}
  \arbvaleurs
  \arbre{A}{B}{0{,}3}{\dots}{0{,}6}{\dots}{\dots}{0{,}25}
\end{tikzpicture}
\end{minipage}

\etape{À vous de jouer}
\begin{minipage}[c]{0.44\linewidth}\raggedright
On donne l'arbre pondéré ci-contre.
\begin{questions}[label=\alph*)]
  \item Écrire chaque nombre donné sous la forme d'une probabilité.
  \item Compléter l'arbre.
  \item Calculer $P(A\cap\overline{B})$ et $P(\overline{A}\cap B)$.
\end{questions}
\end{minipage}\hfill
\begin{minipage}[c]{0.54\linewidth}\centering
\begin{tikzpicture}
  \arbvaleurs
  \arbre{A}{B}{0{,}8}{\dots}{\dots}{0{,}1}{0{,}35}{\dots}
\end{tikzpicture}
\end{minipage}

\end{methodebox}

\begin{methodebox}[Méthode 6 --- Construire un arbre à partir d'un énoncé]
\etapeexemples
Au self du lycée, $60\,\%$ des élèves choisissent le plat végétarien, les autres le
plat avec viande. Parmi les élèves qui prennent le plat végétarien, $30\,\%$ prennent
une salade de fruits en dessert ; parmi les autres, $45\,\%$. On choisit un élève au
hasard. On note $V$ : \og il prend le plat végétarien\fg{} et $F$ : \og il prend une
salade de fruits\fg.
\begin{questions}[label=\alph*)]
  \item Construire l'arbre pondéré qui décrit la situation.
  \item Calculer $P(V\cap F)$ et interpréter le résultat.
  \item Calculer $P(\overline{V}\cap\overline{F})$ et interpréter le résultat.
\end{questions}

\etape{À vous de jouer}
Dans un club de football, $25\,\%$ des licenciés sont des filles. Parmi les filles,
$40\,\%$ jouent en compétition ; parmi les garçons, $60\,\%$. On choisit un licencié au
hasard. On note $F$ : \og c'est une fille\fg{} et $C$ : \og il joue en compétition\fg.
\begin{questions}[label=\alph*)]
  \item Construire l'arbre pondéré qui décrit la situation.
  \item Calculer $P(F\cap C)$ et interpréter le résultat.
  \item Calculer $P(\overline{F}\cap\overline{C})$ et interpréter le résultat.
\end{questions}

\end{methodebox}

\afaire{exercices 13 à 16}

\partiecours{Formule des probabilités totales}

\begin{proprietebox}[Propriété 4 --- Formule des probabilités totales]
\begin{minipage}[c]{0.44\linewidth}\raggedright
La probabilité d'un événement $B$ est la \textbf{somme des probabilités des chemins}
qui mènent à $B$ :
\[
  P(B)=P(A\cap B)+P(\overline{A}\cap B).
\]
Plus généralement, si $A_{1}$, $A_{2}$, $A_{3}$ sont incompatibles deux à deux et
couvrent tout $\Omega$ (chaque issue est dans un et un seul), alors
\[
  P(B)=P(A_{1}\cap B)+P(A_{2}\cap B)+P(A_{3}\cap B).
\]
\end{minipage}\hfill
\begin{minipage}[c]{0.54\linewidth}\centering
\begin{tikzpicture}[arbA/.style={chemin},arbAB/.style={chemin},arbnA/.style={chemin},
    arbnAB/.style={chemin}]
  \arbsansfeuilles
  \arbre{A}{B}{}{}{}{}{}{}
  \feuille{AB}{\color{coulExo}$P(A\cap B)$}
  \feuille{AnB}{$P(A\cap\overline{B})$}
  \feuille{nAB}{\color{coulExo}$P(\overline{A}\cap B)$}
  \feuille{nAnB}{$P(\overline{A}\cap\overline{B})$}
\end{tikzpicture}
\end{minipage}
\end{proprietebox}

\begin{methodebox}[Méthode 7 --- Calculer $P(B)$, puis une probabilité \og inversée\fg]
\etapeexemples
\textbf{Exemple 1.} Une messagerie reçoit $20\,\%$ de messages indésirables
(\og spams\fg). Son filtre envoie dans le dossier \og Indésirables\fg{} $90\,\%$ des
spams et $5\,\%$ des autres messages. On choisit un message au hasard. On note $S$ :
\og c'est un spam\fg{} et $I$ : \og il est envoyé dans les Indésirables\fg.
\begin{questions}[label=\alph*)]
  \item Construire l'arbre pondéré.
  \item Montrer que $P(I)=0{,}22$.
  \item Un message est dans les Indésirables. Calculer la probabilité que ce soit un
    spam (fraction irréductible), puis interpréter.
\end{questions}

\medskip
\textbf{Exemple 2.} Trois machines fabriquent des pièces : la machine $M_{1}$ en
fabrique $50\,\%$, $M_{2}$ en fabrique $30\,\%$ et $M_{3}$ le reste. Une pièce est
défectueuse dans $2\,\%$ des cas pour $M_{1}$, $3\,\%$ pour $M_{2}$ et $5\,\%$ pour
$M_{3}$. On choisit une pièce au hasard ; $D$ : \og la pièce est défectueuse\fg.
\begin{questions}[label=\alph*)]
  \item Construire l'arbre pondéré.
  \item Calculer $P(D)$.
\end{questions}

\etape{À vous de jouer}
\textbf{1.} Dans un groupe, $10\,\%$ des personnes sont malades. Un test rapide est
positif pour $80\,\%$ des malades et pour $5\,\%$ des personnes qui ne sont pas
malades. On choisit une personne au hasard ; $M$ : \og elle est malade\fg{} et
$T$ : \og son test est positif\fg.
\begin{questions}[label=\alph*)]
  \item Construire l'arbre pondéré.
  \item Montrer que $P(T)=0{,}125$.
  \item Le test est positif. Calculer la probabilité que la personne soit malade
    (fraction irréductible), puis interpréter.
\end{questions}
\textbf{2.} Un magasin en ligne confie $60\,\%$ de ses colis au transporteur $T_{1}$,
$30\,\%$ à $T_{2}$ et le reste à $T_{3}$. Un colis arrive en retard dans $5\,\%$ des cas
avec $T_{1}$, $10\,\%$ avec $T_{2}$ et $20\,\%$ avec $T_{3}$. On choisit un colis au
hasard ; $R$ : \og il arrive en retard\fg.
\begin{questions}[label=\alph*)]
  \item Construire l'arbre pondéré.
  \item Calculer $P(R)$.
\end{questions}

\end{methodebox}

\begin{methodebox}[Méthode 8 --- Retrouver une probabilité inconnue $x$]
\etapeexemples
Dans un club de randonnée, $40\,\%$ des adhérents ont plus de $50$ ans. Parmi eux,
$70\,\%$ participent à la sortie du dimanche. Au total, $52\,\%$ des adhérents
participent à la sortie du dimanche. On choisit un adhérent au hasard ; $A$ : \og il a
plus de $50$ ans\fg{} et $D$ : \og il participe à la sortie\fg. On note
$x=P_{\overline{A}}(D)$.
\begin{questions}[label=\alph*)]
  \item Construire l'arbre pondéré, avec $x$.
  \item Montrer que $0{,}6x+0{,}28=0{,}52$.
  \item Résoudre cette équation et interpréter la solution.
\end{questions}

\etape{À vous de jouer}
Dans une médiathèque, $25\,\%$ des inscrits ont moins de $18$ ans. Parmi eux, $60\,\%$
empruntent des bandes dessinées. Au total, $51\,\%$ des inscrits empruntent des bandes
dessinées. On choisit un inscrit au hasard ; $A$ : \og il a moins de $18$ ans\fg{} et
$D$ : \og il emprunte des bandes dessinées\fg. On note $x=P_{\overline{A}}(D)$.
\begin{questions}[label=\alph*)]
  \item Construire l'arbre pondéré, avec $x$.
  \item Montrer que $0{,}75x+0{,}15=0{,}51$.
  \item Résoudre cette équation et interpréter la solution.
\end{questions}

\end{methodebox}

\afaire{exercices 17 à 20}

\partiecours{Événements indépendants}

\begin{definitionbox}[Définition 4 --- Événements indépendants]
$B$ est \textbf{indépendant} de $A$ quand savoir que $A$ est réalisé ne change pas la
probabilité de $B$ : $P_{A}(B)=P(B)$ (avec $P(A)\neq 0$). On remplace alors $P_{A}(B)$
dans la formule des probabilités conditionnelles :
\begin{calculs}
P(A\cap B) &= P(A)\times P_{A}(B) &\qquad& \rem{propriété 2}\\
P(A\cap B) &= P(A)\times P(B) && \rem{car $P_{A}(B)=P(B)$}
\end{calculs}
D'où la définition : deux événements $A$ et $B$ sont \textbf{indépendants} quand
\[
  P(A\cap B)=P(A)\times P(B).
\]
\end{definitionbox}

\begin{proprietebox}[Propriété 5 --- Contraires d'événements indépendants]
Si $A$ et $B$ sont indépendants, alors $\overline{A}$ et $B$ le sont aussi, de même que
$A$ et $\overline{B}$, et $\overline{A}$ et $\overline{B}$.
\end{proprietebox}

\begin{remarquebox}[Indépendants ou incompatibles ?]
Ce sont deux notions différentes. Incompatibles : $A$ et $B$ ne peuvent pas se produire
ensemble, $P(A\cap B)=0$. Indépendants : l'un n'a pas d'influence sur l'autre,
$P(A\cap B)=P(A)\times P(B)$.
\end{remarquebox}

\begin{methodebox}[Méthode 9 --- Démontrer que deux événements sont, ou non, indépendants]
\etapeexemples
\begin{minipage}[c]{0.44\linewidth}\raggedright
\begin{questions}[label=\alph*)]
  \item On lance un dé équilibré à $12$ faces, numérotées de $1$ à $12$. On note
    $A$ : \og le nombre est pair\fg{} et $B$ : \og le nombre est un multiple de
    $3$\fg. $A$ et $B$ sont-ils indépendants ?
  \item On donne l'arbre ci-contre. $A$ et $B$ sont-ils indépendants ?
\end{questions}
\end{minipage}\hfill
\begin{minipage}[c]{0.54\linewidth}\centering
\begin{tikzpicture}
  \arbre{A}{B}{0{,}4}{0{,}6}{0{,}3}{0{,}7}{0{,}5}{0{,}5}
\end{tikzpicture}
\end{minipage}

\etape{À vous de jouer}
\begin{minipage}[c]{0.44\linewidth}\raggedright
\begin{questions}[label=\alph*)]
  \item On lance un dé équilibré à $10$ faces, numérotées de $1$ à $10$. On note
    $A$ : \og le nombre est pair\fg{} et $B$ : \og le nombre est inférieur ou égal à
    $4$\fg. $A$ et $B$ sont-ils indépendants ?
  \item On donne l'arbre ci-contre. $A$ et $B$ sont-ils indépendants ?
\end{questions}
\end{minipage}\hfill
\begin{minipage}[c]{0.54\linewidth}\centering
\begin{tikzpicture}
  \arbre{A}{B}{0{,}6}{0{,}4}{0{,}25}{0{,}75}{0{,}25}{0{,}75}
\end{tikzpicture}
\end{minipage}

\end{methodebox}

\begin{methodebox}[Méthode 10 --- Calculer avec deux événements indépendants]
\etapeexemples
Une maison a deux détecteurs de fumée qui fonctionnent indépendamment l'un de l'autre.
Dans l'année, le premier tombe en panne avec une probabilité de $0{,}05$, le second
avec une probabilité de $0{,}08$. On note $A$ et $B$ ces deux événements.
\begin{questions}[label=\alph*)]
  \item Calculer la probabilité que les deux détecteurs tombent en panne.
  \item Calculer la probabilité qu'au moins un détecteur tombe en panne.
  \item Calculer la probabilité qu'aucun ne tombe en panne.
\end{questions}

\etape{À vous de jouer}
Pour aller au lycée, Inès prend un bus puis un train, dont les retards sont
indépendants. Le bus est en retard avec une probabilité de $0{,}1$, le train avec une
probabilité de $0{,}2$. On note $A$ et $B$ ces deux événements.
\begin{questions}[label=\alph*)]
  \item Calculer la probabilité que le bus et le train soient en retard.
  \item Calculer la probabilité qu'au moins l'un des deux soit en retard.
  \item Calculer la probabilité qu'aucun ne soit en retard.
\end{questions}

\end{methodebox}

\afaire{exercices 21 à 24}

\partiecours{Trois façons de tirer deux objets}

\begin{definitionbox}[Définition 5 --- Tirages avec remise, sans remise, simultanés]
On tire deux boules dans une urne.
\begin{itemize}
  \item \textbf{Successifs avec remise} : on remet la première boule avant de tirer la
    seconde. L'urne est la même aux deux tirages : les deux tirages sont
    \textbf{indépendants}, et l'arbre porte les mêmes probabilités aux deux niveaux.
  \item \textbf{Successifs sans remise} : la première boule reste dehors. Au second
    tirage, il manque une boule : les probabilités du second niveau changent selon la
    première boule (probabilités conditionnelles).
  \item \textbf{Simultanés} : on prend les deux boules d'un coup, sans ordre. On liste
    les paires possibles, toutes équiprobables ; on obtient les mêmes probabilités que
    sans remise.
\end{itemize}
\end{definitionbox}

\begin{methodebox}[Méthode 11 --- Calculer une probabilité dans un tirage de deux boules]
\etapeexemples
Une urne contient $3$ boules rouges et $2$ boules vertes. On tire deux boules au hasard
et on note $M$ : \og les deux boules ont la même couleur\fg. Calculer $P(M)$ quand :
\begin{questions}[label=\alph*)]
  \item on tire une boule, on la remet, puis on en tire une seconde ;
  \item on tire une boule, on la garde, puis on en tire une seconde ;
  \item on tire les deux boules en même temps.
\end{questions}

\etape{À vous de jouer}
Un sac contient $4$ jetons bleus et $2$ jetons jaunes. On tire deux jetons au hasard et
on note $M$ : \og les deux jetons ont la même couleur\fg. Calculer $P(M)$ quand :
\begin{questions}[label=\alph*)]
  \item on tire un jeton, on le remet, puis on en tire un second ;
  \item on tire un jeton, on le garde, puis on en tire un second ;
  \item on tire les deux jetons en même temps.
\end{questions}

\end{methodebox}

\afaire{exercices 25 à 27}

\partiecours{Épreuve de Bernoulli}

\begin{definitionbox}[Définition 6 --- Épreuve de Bernoulli]
Une \textbf{épreuve de Bernoulli} est une expérience aléatoire à deux issues :
\textbf{succès} $S$, de probabilité $p$, et \textbf{échec} $\overline{S}$, de probabilité
$1-p$. Le nombre $p$ est le \textbf{paramètre} de l'épreuve.

\smallskip
Exemples : obtenir \og pile\fg{} avec une pièce équilibrée ($p=0{,}5$) ; obtenir $6$
avec un dé équilibré ($p=\frac{1}{6}$) ; réussir un tir au but.
\end{definitionbox}

\begin{proprietebox}[Propriété 6 --- Répéter une épreuve de Bernoulli]
Quand on répète la même épreuve de Bernoulli de façon \textbf{indépendante}, l'arbre
porte $p$ sur chaque branche $S$ et $1-p$ sur chaque branche $\overline{S}$ ; la
probabilité d'un chemin est le produit de ses branches.
\end{proprietebox}

\begin{methodebox}[Méthode 12 --- Répéter trois fois une épreuve de Bernoulli]
\etapeexemples
Une joueuse de basket réussit un lancer franc avec une probabilité de $0{,}8$. Elle
tire trois lancers francs, de façon indépendante ; $S$ : \og le lancer est réussi\fg.
\begin{questions}[label=\alph*)]
  \item Construire l'arbre pondéré.
  \item Calculer la probabilité qu'elle réussisse les trois lancers.
  \item Calculer la probabilité qu'elle en réussisse exactement deux.
  \item Calculer la probabilité qu'elle en réussisse au moins un.
\end{questions}

\etape{À vous de jouer}
Une graine germe avec une probabilité de $0{,}9$. On plante trois graines, qui germent
indépendamment les unes des autres ; $S$ : \og la graine germe\fg.
\begin{questions}[label=\alph*)]
  \item Construire l'arbre pondéré.
  \item Calculer la probabilité que les trois graines germent.
  \item Calculer la probabilité qu'exactement deux graines germent.
  \item Calculer la probabilité qu'au moins une graine germe.
\end{questions}

\end{methodebox}

\afaire{exercices 28 à 30}

\partiecours{Variable aléatoire}

\begin{definitionbox}[Définition 7 --- Variable aléatoire]
Une \textbf{variable aléatoire} $X$ associe un nombre réel à chaque issue d'une
expérience aléatoire (un gain, un nombre de succès\dots).

\smallskip
$\{X=a\}$ est l'événement \og $X$ prend la valeur $a$\fg ; sa probabilité est notée
$P(X=a)$. De même, $\{X\ge a\}$ est l'événement \og $X$ prend une valeur supérieure ou
égale à $a$\fg.
\end{definitionbox}

\begin{remarquebox}[Une \og variable\fg{} qui est en fait une fonction]
Le mot est trompeur : dans $f(x)$, $x$ est une variable ; ici, $X$ n'est pas une
variable mais une \textbf{fonction}, qui à chaque issue associe un nombre :
\[
  \text{issue}\longmapsto X(\text{issue}).
\]
Sa valeur dépend du hasard, puisque l'issue n'est connue qu'après l'expérience.
\end{remarquebox}

\textbf{Exemple 1.} On lance deux fois une pièce ($P$ : pile, $F$ : face) et on note
$X$ le nombre de \og pile\fg{} obtenus.
\begin{center}
\renewcommand{\arraystretch}{1.4}
\begin{tabular}{|c|c|c|c|c|}
\hline
Issue & $PP$ & $PF$ & $FP$ & $FF$\\\hline
Valeur de $X$ & $2$ & $1$ & $1$ & $0$\\\hline
\end{tabular}
\end{center}
L'événement $\{X=1\}$ est $\ens{PF\,;FP}$ : $2$ issues sur $4$, donc $P(X=1)=\frac{1}{2}$.

\medskip
\textbf{Exemple 2.} On lance un dé : on gagne $3$\,€ si on obtient $6$, sinon on perd
$1$\,€. On note $X$ le gain, en euros.
\begin{center}
\renewcommand{\arraystretch}{1.4}
\begin{tabular}{|c|c|c|c|c|c|c|}
\hline
Issue & $1$ & $2$ & $3$ & $4$ & $5$ & $6$\\\hline
Valeur de $X$ & $-1$ & $-1$ & $-1$ & $-1$ & $-1$ & $3$\\\hline
\end{tabular}
\end{center}
L'événement $\{X=-1\}$ est $\ens{1\,;2\,;3\,;4\,;5}$, donc $P(X=-1)=\frac{5}{6}$.

\begin{methodebox}[Méthode 13 --- Calculer des probabilités avec une variable aléatoire]
\etapeexemples
On lance un dé équilibré à six faces. On gagne $5$\,€ si on obtient $6$, on gagne
$1$\,€ si on obtient $4$ ou $5$, et on perd $2$\,€ sinon. On note $X$ le gain, en euros
(une perte est un gain négatif).
\begin{questions}[label=\alph*)]
  \item Donner les valeurs prises par $X$.
  \item Calculer $P(X=1)$ et $P(X=-2)$.
  \item Calculer $P(X\ge 1)$.
  \item Décrire par une phrase l'événement $\{X\le 1\}$ et calculer sa probabilité.
\end{questions}

\etape{À vous de jouer}
Une roue est partagée en $8$ secteurs égaux : $1$ secteur fait gagner $10$\,€, $3$
secteurs font gagner $2$\,€ et les $4$ autres font perdre $1$\,€. On fait tourner la
roue et on note $X$ le gain, en euros.
\begin{questions}[label=\alph*)]
  \item Donner les valeurs prises par $X$.
  \item Calculer $P(X=2)$ et $P(X=-1)$.
  \item Calculer $P(X\ge 2)$.
  \item Décrire par une phrase l'événement $\{X\le 2\}$ et calculer sa probabilité.
\end{questions}

\end{methodebox}

\afaire{exercices 31 à 33}

\partiecours{Loi de probabilité}

\begin{definitionbox}[Définition 8 --- Loi de probabilité d'une variable aléatoire]
La \textbf{loi de probabilité} d'une variable aléatoire $X$ est en fait un
\textbf{tableau} : sur la première ligne, les valeurs que $X$ peut prendre ; sous chacune,
sa probabilité d'apparition.

\smallskip
Si $X$ prend les valeurs $x_{1}$, $x_{2}$, \dots, $x_{n}$, on note
$p_{i}=P(X=x_{i})$ :
\begin{center}
\renewcommand{\arraystretch}{1.4}
\begin{tabular}{|c|c|c|c|c|}
\hline
$x_{i}$ & $x_{1}$ & $x_{2}$ & \dots & $x_{n}$\\\hline
$P(X=x_{i})$ & $p_{1}$ & $p_{2}$ & \dots & $p_{n}$\\\hline
\end{tabular}
\end{center}
La somme des probabilités vaut $1$ : $p_{1}+p_{2}+\dots+p_{n}=1$.

\smallskip
Exemple : nombre $X$ de \og pile\fg{} en deux lancers (partie précédente) ;
$0$ pile pour $FF$, $1$ pour $PF$ et $FP$, $2$ pour $PP$.
\begin{center}
\renewcommand{\arraystretch}{2.4}
\begin{tabular}{|c|c|c|c|}
\hline
$x_{i}$ & $0$ & $1$ & $2$\\\hline
$P(X=x_{i})$ & $\dfrac{1}{4}$ & $\dfrac{2}{4}$ & $\dfrac{1}{4}$\\\hline
\end{tabular}
\end{center}

\smallskip
\textbf{Loi de Bernoulli de paramètre $p$} : $X$ vaut $1$ en cas de succès et $0$ en cas
d'échec ; ${P(X=1)=p}$ et ${P(X=0)=1-p}$.
\end{definitionbox}

\begin{methodebox}[Méthode 14 --- Déterminer la loi d'une variable aléatoire]
\etapeexemples
\textbf{Exemple 1.} Pour jouer à un jeu en ligne, on paie $2$\,€. Un sac virtuel
contient $3$ jetons gagnants sur $10$. On tire un jeton, on le remet, puis on en tire
un second. Chaque jeton gagnant rapporte $5$\,€. On note $X$ le gain algébrique : ce
que l'on reçoit, moins les $2$\,€ payés.
\begin{questions}[label=\alph*)]
  \item Donner les valeurs prises par $X$.
  \item Déterminer la loi de probabilité de $X$.
\end{questions}

\medskip
\textbf{Exemple 2.} La loi d'une variable aléatoire $Y$ est donnée par le tableau
ci-dessous. Calculer $a$, puis $P(Y\ge 1)$.
\begin{center}
\renewcommand{\arraystretch}{1.4}
\begin{tabular}{|c|c|c|c|c|}
\hline
$y_{i}$ & $0$ & $1$ & $2$ & $3$\\\hline
$P(Y=y_{i})$ & $0{,}1$ & $0{,}3$ & $a$ & $0{,}2$\\\hline
\end{tabular}
\end{center}

\etape{À vous de jouer}
\textbf{1.} Pour jouer, on paie $5$\,€. Un sac contient $4$ jetons gagnants sur $10$.
On tire un jeton, on le remet, puis on en tire un second. Chaque jeton gagnant rapporte
$5$\,€. On note $X$ le gain algébrique.
\begin{questions}[label=\alph*)]
  \item Donner les valeurs prises par $X$.
  \item Déterminer la loi de probabilité de $X$.
\end{questions}
\textbf{2.} La loi d'une variable aléatoire $Y$ est donnée ci-dessous. Calculer $a$, puis
$P(Y\ge 0)$.
\begin{center}
\renewcommand{\arraystretch}{1.4}
\begin{tabular}{|c|c|c|c|c|}
\hline
$y_{i}$ & $-1$ & $0$ & $2$ & $5$\\\hline
$P(Y=y_{i})$ & $0{,}25$ & $a$ & $0{,}35$ & $0{,}1$\\\hline
\end{tabular}
\end{center}

\end{methodebox}

\afaire{exercices 34 à 37}

\partiecours{Espérance, variance, écart-type}

\begin{definitionbox}[Définition 9 --- Espérance, variance, écart-type]
Soit $X$ une variable aléatoire de loi $p_{i}=P(X=x_{i})$, pour $i$ de $1$ à $n$.
\begin{itemize}
  \item L'\textbf{espérance} de $X$ : $E(X)=p_{1}x_{1}+p_{2}x_{2}+\dots+p_{n}x_{n}$,
    la somme des valeurs $x_{i}$ multipliées par leurs probabilités.
  \item La \textbf{variance} de $X$ :
    $V(X)=p_{1}\big(x_{1}-E(X)\big)^{2}+p_{2}\big(x_{2}-E(X)\big)^{2}+\dots
    +p_{n}\big(x_{n}-E(X)\big)^{2}$, la somme des carrés des écarts $x_{i}-E(X)$
    multipliés par leurs probabilités.
  \item L'\textbf{écart-type} de $X$ : $\sigma(X)=\sqrt{V(X)}$.
\end{itemize}
\end{definitionbox}

\begin{remarquebox}[Ce que disent ces nombres]
\begin{itemize}
  \item $E(X)$ est la \textbf{moyenne} des valeurs de $X$ quand on répète l'expérience
    un grand nombre de fois. Pour un gain : jeu \textbf{favorable} au joueur si
    $E(X)>0$, \textbf{équitable} si $E(X)=0$, \textbf{défavorable} si $E(X)<0$.
  \item $\sigma(X)$ mesure la \textbf{dispersion} des valeurs autour de $E(X)$ : plus
    il est grand, plus les résultats sont variables (plus le jeu est risqué).
  \item Loi de Bernoulli de paramètre $p$ : $E(X)=p\times 1+(1-p)\times 0=p$.
  \item Calculatrice : en statistiques, les $x_{i}$ en valeurs et les $p_{i}$ en
    effectifs ; la moyenne affichée est $E(X)$ et l'écart-type $\sigma(X)$.
\end{itemize}
\end{remarquebox}

\begin{methodebox}[Méthode 15 --- Calculer et interpréter une espérance]
\etapeexemples
On reprend le jeu de l'exemple 1 de la méthode 14 (mise de $2$\,€, chaque jeton gagnant
rapporte $5$\,€) :
\begin{center}
\renewcommand{\arraystretch}{1.4}
\begin{tabular}{|c|c|c|c|}
\hline
$x_{i}$ & $-2$ & $3$ & $8$\\\hline
$P(X=x_{i})$ & $0{,}49$ & $0{,}42$ & $0{,}09$\\\hline
\end{tabular}
\end{center}
\begin{questions}[label=\alph*)]
  \item Calculer $E(X)$.
  \item Interpréter le résultat. Le jeu est-il favorable au joueur ?
  \item L'organisateur remplace la mise de $2$\,€ par une mise de $m$ euros. Quelle mise
    rend le jeu équitable ?
\end{questions}

\etape{À vous de jouer}
On reprend le jeu de l'exercice 1 de la méthode 14 (mise de $5$\,€, chaque jeton
gagnant rapporte $5$\,€) :
\begin{center}
\renewcommand{\arraystretch}{1.4}
\begin{tabular}{|c|c|c|c|}
\hline
$x_{i}$ & $-5$ & $0$ & $5$\\\hline
$P(X=x_{i})$ & $0{,}36$ & $0{,}48$ & $0{,}16$\\\hline
\end{tabular}
\end{center}
\begin{questions}[label=\alph*)]
  \item Calculer $E(X)$.
  \item Interpréter le résultat. Le jeu est-il favorable au joueur ?
  \item L'organisateur remplace la mise de $5$\,€ par une mise de $m$ euros. Quelle mise
    rend le jeu équitable ?
\end{questions}

\end{methodebox}

\begin{methodebox}[Méthode 16 --- Calculer une variance et un écart-type]
\etapeexemples
Pour le jeu de la méthode 15 (mise de $2$\,€), on a trouvé $E(X)=1$. Calculer $V(X)$,
puis $\sigma(X)$ arrondi au centième.

\etape{À vous de jouer}
Pour le jeu de l'« À vous de jouer » de la méthode 15 (mise de $5$\,€), on a trouvé
$E(X)=-1$. Calculer $V(X)$, puis $\sigma(X)$ arrondi au centième.

\end{methodebox}
\afaire{exercices 38 à 45}

\end{document}
